\documentclass[12pt]{article}

% -------------------------
% PAQUETES NECESARIOS
% -------------------------
\usepackage[utf8]{inputenc}
\usepackage[T1]{fontenc}
\usepackage{titlesec}
\usepackage{fancyhdr}
\usepackage{graphicx}
\usepackage{listings}
\usepackage{xcolor}
\usepackage{geometry}
\usepackage[spanish]{babel}
\usepackage{graphicx}
\usepackage{booktabs} % Para líneas profesionales (\toprule, \midrule, \bottomrule)
\usepackage{multirow} % Para agrupar celdas verticales

\usepackage{amsmath}
\usepackage{booktabs}
\usepackage{xcolor}
\usepackage{pgfplots}
\pgfplotsset{compat=1.18}
\usepackage[colorlinks=true, linkcolor=blue, citecolor=blue, urlcolor=blue]{hyperref}


\usepackage{algorithm}
\usepackage{algpseudocode}
\floatname{algorithm}{Algoritmo}
\algrenewcommand\algorithmicif{\textbf{si}}
\algrenewcommand\algorithmicthen{\textbf{entonces}}
\algrenewcommand\algorithmicelse{\textbf{si no}}
\algrenewcommand\algorithmicend{\textbf{fin}}
\algrenewcommand\algorithmicfor{\textbf{para}}
\algrenewcommand\algorithmicdo{\textbf{hacer}}
\algrenewcommand\algorithmicwhile{\textbf{mientras}}
\algrenewcommand\algorithmicfunction{\textbf{función}}
\algrenewcommand\algorithmicreturn{\textbf{devolver}}
\algrenewtext{EndIf}{\textbf{fin si}}
\algrenewtext{EndFor}{\textbf{fin para}}
\algrenewtext{EndWhile}{\textbf{fin mientras}}
\algrenewtext{EndFunction}{\textbf{fin función}}



% Definición de colores principales
\definecolor{primary}{RGB}{30, 64, 175}
\definecolor{secondary}{RGB}{13, 148, 136}
\definecolor{danger}{RGB}{225, 29, 72}
\definecolor{neutral-bg}{RGB}{248, 250, 252}


% -------------------------
% CONFIGURACIÓN DE PÁGINA
% -------------------------
\geometry{margin=2.5cm}

% Fuente tipo Times New Roman equivalente
\usepackage{mathptmx}

% Interlineado sencillo (por defecto en LaTeX)

% -------------------------
% CONFIGURACIÓN DE ENCABEZADO Y PIE
% -------------------------
\pagestyle{fancy}
\fancyhf{}

\fancyhead[L]{Metodología de la programación paralela}
\fancyhead[R]{2026/2027}

\fancyfoot[L]{Andrey, Abdellah} % <-- EDITAR
\fancyfoot[R]{\thepage}

\renewcommand{\headrulewidth}{0.4pt}
\renewcommand{\footrulewidth}{0.4pt}


\usepackage{listings}
\usepackage{xcolor}

% --- Colores generales ---
\definecolor{mGreen}{rgb}{0,0.6,0}
\definecolor{mGray}{rgb}{0.5,0.5,0.5}
\definecolor{mPurple}{rgb}{0.58,0,0.82}
\definecolor{bgC}{rgb}{0.98,0.98,0.98}
\definecolor{bgValgrind}{rgb}{0.93,0.93,0.93}

% --- ESTILO 1: Para Código C ---
\lstdefinestyle{EstiloC}{
    backgroundcolor=\color{bgC},
    commentstyle=\color{mGreen}\itshape,
    keywordstyle=\color{magenta}\bfseries,
    numberstyle=\tiny\color{mGray},
    stringstyle=\color{mPurple},
    basicstyle=\ttfamily\footnotesize,
    breakatwhitespace=false,
    breaklines=true,                 % ¡Evita que el código se salga por la derecha!
    captionpos=b,
    keepspaces=true,
    numbers=left,
    numbersep=8pt,
    showstringspaces=false,
    tabsize=4,
    language=C,
    frame=single,
    rulecolor=\color{mGray!30}
}

% --- ESTILO 2: Para trazas/logs de Valgrind ---
\lstdefinestyle{EstiloValgrind}{
    backgroundcolor=\color{bgValgrind},
    basicstyle=\ttfamily\scriptsize, % Letra ligeramente más pequeña para trazas anchas
    breakatwhitespace=false,
    breaklines=true,                 % ¡Evita que los reportes de Valgrind se corten!
    captionpos=b,
    frame=single,
    numbers=none,                    % Sin números de línea para logs de terminal
    showstringspaces=false,
    language={},                     % Sin resaltado de sintaxis de lenguaje
    rulecolor=\color{mGray!50}
}

% -------------------------
% INICIO DEL DOCUMENTO
% -------------------------
% acentos y eñes dentro de los listados de código (sin esto, listings falla con UTF-8)
\lstset{literate=
  {á}{{\'a}}1 {é}{{\'e}}1 {í}{{\'i}}1 {ó}{{\'o}}1 {ú}{{\'u}}1
  {Á}{{\'A}}1 {É}{{\'E}}1 {Í}{{\'I}}1 {Ó}{{\'O}}1 {Ú}{{\'U}}1
  {ñ}{{\~n}}1 {Ñ}{{\~N}}1 {ü}{{\"u}}1}

\begin{document}

% =========================
% PORTADA
% =========================
\thispagestyle{empty}

% Bloque arriba a la derecha
\begin{flushright}
Abdellah Bensaidi, Souri 1.3\\
Morales Vicuña, Andrey 1.1\\
Asignatura: Metodología de la programación paralela \\
Grupo: 1 \\
Práctica 3: Paralelización híbrida con MPI+OpenMP. Modelado y auto-optimización \\
\end{flushright}

\vfill

\begin{center}
\includegraphics[width=0.4\textwidth]{sello.png}
\end{center}

\vfill

\begin{center}
\today
\end{center}

\newpage
\tableofcontents
\newpage

% =========================
% SISTEMA COMPUTACIONAL
% =========================
\section{Sistema computacional}

% >>> DATO: generar en el ordenador de las ejecuciones finales con
%   lstopo -f -p --no-legend --no-io sistema.svg
% y convertirlo para LaTeX con:  rsvg-convert -f pdf -o sistema.pdf sistema.svg
% (o directamente: lstopo -f -p --no-legend --no-io sistema.pdf)
\begin{figure}[htbp]
    \centering
    \includegraphics[width=0.8\textwidth]{sistema.pdf}
    \caption{Sistema utilizado (salida de \texttt{lstopo -f -p -{}-no-legend -{}-no-io}).}
    \label{fig:sistema}
\end{figure}

% >>> DATO: descripción de la máquina de las ejecuciones finales
Todas las ejecuciones se han realizado en el mismo sistema: un procesador Intel Core i7-1355U (13.ª generación) con 10 núcleos físicos y 12 hilos hardware. Es un procesador híbrido: tiene 2 núcleos de rendimiento con dos hilos cada uno (\textit{Hyper-Threading}) y 8 núcleos de eficiencia con un hilo cada uno, más lentos. La frecuencia máxima es de 5.0 GHz. Se ha usado Ubuntu, \texttt{gcc} con \texttt{-O3} y OpenMPI.

\newpage

\section{Ejercicios}

\begin{enumerate}

% =====================================================================
%  PARTE 1
% =====================================================================
\item[] \textbf{\large Parte 1. Implementación de paralelismo híbrido (MPI + OpenMP)}

\item \textbf{Cuestión 1.} Con el fin de obtener un código más eficiente en la parte de memoria compartida y poder maximizar el proceso de mejora de la población de individuos: reemplazar las llamadas a la función \texttt{rand} por \texttt{rand\_r} (que es ``thread safe'') para evitar, por un lado, que se genere la misma secuencia de valores aleatorios en diferentes hilos y, por otro, que se pierda paralelismo durante la generación de dichos valores debido a la secuencialidad que introduce la función \texttt{rand}.

% --- Respuesta de la cuestión 1 ---

Se parte del código de la cuestión 6 de la P2 (modelo de islas con comunicaciones colectivas), que es el que se usará en la versión híbrida. Todos los números aleatorios del algoritmo se obtienen con la función \texttt{aleatorio}, así que basta con cambiarla a ella y a la inicialización de la semilla.

\medskip
\noindent\textbf{Generación de números aleatorios (\texttt{mh.c})}\par
\smallskip

\begin{lstlisting}[style=EstiloC]
// rand_r recibe el estado (la semilla) como parámetro, así que es "thread safe", y con
// threadprivate cada hilo tiene su propia copia de la semilla
static unsigned int semilla = 1;
#pragma omp threadprivate(semilla)

int aleatorio(int n) {
	return (rand_r(&semilla) % n);  // genera un numero aleatorio entre 0 y n-1
}
\end{lstlisting}

\textbf{Justificación:} \texttt{rand} guarda su estado en una variable global que comparten todos los hilos. Si varios hilos la llaman a la vez, se pueden pisar el estado y repetir números, y además la biblioteca la protege con un cerrojo, así que los hilos acaban generando los números de uno en uno. \texttt{rand\_r} recibe el estado como parámetro, y con \texttt{threadprivate} cada hilo tiene su propia semilla, de modo que no se comparte nada y cada hilo genera su propia secuencia. Se ha hecho igual que en la P1.

\vspace{0.5cm}

\medskip
\noindent\textbf{Inicialización de la semilla (\texttt{aplicar\_mh} en \texttt{mh.c})}\par
\smallskip

\begin{lstlisting}[style=EstiloC]
// P3 CUESTIÓN 1: en lugar de srand se inicializa la semilla de rand_r
semilla = (unsigned int) (time(NULL) + getpid() + rank * 1000);
\end{lstlisting}

\textbf{Justificación:} Sustituye al \texttt{srand} de la P2. Se mantiene el \texttt{rank} para que cada isla tenga una semilla distinta y evolucione de forma diferente. En esta cuestión todavía no hay hilos, así que hay una semilla por proceso; en la cuestión 2 se añade el número de hilo para que cada hilo de cada proceso tenga la suya. En el \texttt{Makefile} se ha añadido \texttt{-fopenmp} para que se reconozca la directiva \texttt{threadprivate}. Se ha comprobado que el programa da el mismo fitness que la P2.


\newpage
\item \textbf{Cuestión 2.} Implementar un algoritmo híbrido, tomando como referencia las directivas OpenMP y primitivas MPI que arrojaron mejores resultados en las prácticas 1 y 2. Se deberán utilizar, por tanto, las funciones que se considere más adecuadas para obtener un código optimizado y que permita obtener las mejores prestaciones. Ejecutar el código variando el número de procesos y de hilos, y comparar los resultados obtenidos respecto a la versión secuencial utilizando las métricas de rendimiento de la Tabla 1.

Además, se deberá analizar la mejora obtenida en las prestaciones, considerando solo tiempo de ejecución y tiempo de ejecución + fitness (ver apartado (1) de Observaciones), respecto al algoritmo secuencial y también respecto a la mejor versión paralela obtenida en las prácticas anteriores usando el mismo sistema computacional para todas las versiones.

% --- Respuesta de la cuestión 2 ---

La versión híbrida combina las dos mejores versiones de las prácticas anteriores:

\begin{itemize}
    \item \textbf{Nivel MPI (P2, cuestión 6):} modelo de islas con comunicaciones colectivas. Cada proceso es una isla y cada NGM generaciones el proceso 0 recoge los NEM mejores de cada isla con \texttt{MPI\_Gather} y les devuelve los NEM mejores globales y la decisión de parar con \texttt{MPI\_Bcast}.
    \item \textbf{Nivel OpenMP (P1):} dentro de cada proceso, una única región paralela para todo el algoritmo, con los bucles sobre la población repartidos entre los hilos.
\end{itemize}

A continuación se muestran las partes del código modificadas respecto a la cuestión 1.

% cambio 1: inicialización de MPI con hilos

\medskip
\noindent\textbf{Inicialización de MPI y número de hilos (\texttt{main.c})}\par
\smallskip

\begin{lstlisting}[style=EstiloC]
int provided;
MPI_Init_thread(&argc, &argv, MPI_THREAD_FUNNELED, &provided);
...
if (provided < MPI_THREAD_FUNNELED) {
	...
	MPI_Finalize();
	return EXIT_FAILURE;
}
...
int NHILOS = atoi(argv[8]); // P3: hilos OpenMP que pone en marcha cada proceso
assert(NHILOS >= 1);
omp_set_num_threads(NHILOS);
\end{lstlisting}

\textbf{Justificación:} Cuando un proceso MPI tiene varios hilos hay que avisar a MPI de cómo se van a usar con \texttt{MPI\_Init\_thread}. Se pide el nivel \texttt{MPI\_THREAD\_FUNNELED}, que significa que solo el hilo maestro hace llamadas a MPI. Es el nivel más bajo que nos sirve, porque en nuestro código todas las comunicaciones las hace el maestro (ver más abajo), y los niveles más altos añaden cerrojos dentro de MPI que no necesitamos. Si la biblioteca no lo soporta, el programa termina. El número de hilos por proceso se pasa como un nuevo argumento para poder probar combinaciones de procesos e hilos sin recompilar.

\vspace{0.5cm}

% cambio 2: región paralela

\medskip
\noindent\textbf{Algoritmo híbrido (\texttt{aplicar\_mh} en \texttt{mh.c})}\par
\smallskip

\begin{lstlisting}[style=EstiloC]
MPI_Comm comm = MPI_COMM_WORLD;
MPI_Datatype mpi_int = MPI_INT;

#pragma omp parallel default(none) \
	shared(d, n, m, n_gen, tam_pob, m_rate, sol, NGM, NEM, rank, size, individuo_type, comm, mpi_int, \
	       tam_local, bloque, poblacion, emigrantes, inmigrantes, mezcla, p_mezcla, \
	       fitness_modelo, ngeneracion, continuar, value)
{
	// P3 CUESTIÓN 1: semilla distinta para cada hilo de cada proceso
	int id = 0;
#if defined (_OPENMP)
	id = omp_get_thread_num();
#endif
	semilla = (unsigned int) (time(NULL) + getpid() + rank * 1000 + 7919 * id);

	#pragma omp for schedule(dynamic)
	for (int i = 0; i < tam_local; i++) {
		poblacion[i] = &bloque[i];
		crear_individuo(n, m, poblacion[i]->array_int);
		fitness(d, poblacion[i], n, m);
	}

	#pragma omp single
	{
		qsort(poblacion, tam_local, sizeof(Individuo *), comp_fitness);
		fitness_modelo = poblacion[0]->fitness;
	}

	int g = 0;
	while (g < n_gen && continuar)
	{
		int gen_bloque = 0;
		for (int k = 0; k < NGM && g < n_gen; k++, g++, gen_bloque++)
		{
			#pragma omp for schedule(dynamic)
			for (int i = 0; i < (tam_local/2) - 1; i += 2) {
				cruzar(poblacion[i], poblacion[i+1], poblacion[tam_local/2 + i], poblacion[tam_local/2 + i + 1], n, m);
			}

			#pragma omp for schedule(static)
			for (int i = 0; i < tam_local; i++) {
				if (i >= tam_local/4)
					mutar(poblacion[i], n, m, m_rate);
				fitness(d, poblacion[i], n, m);
			}

			#pragma omp single
			qsort(poblacion, tam_local, sizeof(Individuo *), comp_fitness);
		}

		// Migración: solo el hilo maestro de cada proceso llama a MPI
		#pragma omp master
		{
			copiar_mejores(poblacion, emigrantes, NEM);
			MPI_Gather(emigrantes, NEM, individuo_type, mezcla, NEM, individuo_type, 0, comm);
			if (rank == 0) {
				... // ordenar la mezcla, elegir los NEM mejores y decidir la convergencia
			}
			MPI_Bcast(inmigrantes, NEM, individuo_type, 0, comm);
			MPI_Bcast(&continuar, 1, mpi_int, 0, comm);
			reemplazar_peores(poblacion, tam_local, inmigrantes, NEM);
		}
		#pragma omp barrier
	}

	// Resultado final: el proceso 0 recoge el mejor de cada isla (solo el maestro)
	#pragma omp master
	{
		... // MPI_Gather del mejor de cada isla y copia del mejor global en sol
	}
} // fin de la región paralela (barrera implícita)
\end{lstlisting}

\textbf{Justificación:}

\begin{itemize}
    \item \textbf{Una sola región paralela.} Como en la P1, los hilos se crean una vez al principio y no en cada generación, para no pagar el coste de crearlos muchas veces. Con \texttt{default(none)} estamos obligados a indicar el alcance de cada variable, lo que evita compartir algo sin querer. En OpenMPI, \texttt{MPI\_COMM\_WORLD} y \texttt{MPI\_INT} son variables globales y \texttt{default(none)} da error con ellas, por eso se copian a las variables locales \texttt{comm} y \texttt{mpi\_int}.
    \item \textbf{Semilla por hilo.} A la semilla de la cuestión 1 se le suma el número de hilo multiplicado por un primo, para que cada hilo de cada proceso genere una secuencia distinta.
    \item \textbf{Reparto de los bucles.} Se usan las mismas directivas que dieron mejor resultado en la P1. La población inicial y el cruce usan \texttt{schedule(dynamic)} porque cada iteración tiene un coste distinto (crear un individuo y reparar los duplicados del cruce tardan más o menos según los números que salgan). La mutación y el fitness se han juntado en un solo bucle, porque el fitness de un individuo solo depende de su propia mutación; así hay una barrera menos por generación. Este bucle usa \texttt{schedule(static)}, porque su coste lo marca el fitness, que es igual para todos los individuos. La ordenación con \texttt{qsort} es secuencial y la hace un solo hilo con \texttt{single}. La barrera implícita de cada \texttt{for} es necesaria: la mutación modifica individuos que son padres en el cruce, y no se puede ordenar hasta que todos tengan su fitness.
    \item \textbf{Comunicaciones solo en el hilo maestro.} La migración y la recogida final del resultado van dentro de \texttt{master}, que es lo que permite usar el nivel \texttt{FUNNELED}. Como \texttt{master} no tiene barrera implícita, se pone un \texttt{barrier} después: ningún hilo puede seguir con la población ni leer \texttt{continuar} hasta que el maestro haya terminado la migración. Así todos los hilos de todos los procesos salen del bucle en la misma generación y no hay bloqueos.
    \item \textbf{Variables privadas y compartidas.} El contador de generaciones \texttt{g} es privado, pero todos los hilos hacen las mismas iteraciones, así que vale lo mismo en todos. \texttt{continuar} es compartida: la escribe el maestro y el resto la lee después de la barrera.
\end{itemize}

Se ha comprobado que la solución es correcta volviendo a calcular su fitness a partir de la matriz de distancias. También se ha comprobado que no hay bloqueos con distintas combinaciones de procesos e hilos y con migración en cada generación (NGM = 1).

\vspace{0.5cm}

\medskip
\noindent\textbf{Compilación y ejecución}\par
\smallskip

En el \texttt{Makefile} se ha añadido \texttt{-fopenmp} y el número de hilos como parámetro, y se ejecuta con \texttt{mpirun -{}-bind-to none}. Por defecto OpenMPI fija cada proceso a un núcleo, así que todos los hilos de un proceso acabarían en ese mismo núcleo y no se ganaría nada; con \texttt{-{}-bind-to none} el sistema operativo puede repartir los hilos por todos los núcleos. En una prueba con 1 proceso y 2 hilos, el tiempo bajó de 1.78 s a 0.99 s solo con este cambio.

\vspace{0.5cm}

\medskip
\noindent\textbf{Experimentos}\par
\smallskip

Para comparar todas las versiones en el mismo sistema, como pide el enunciado, se han vuelto a ejecutar en la misma sesión y alternadas, con el script \texttt{experimentos\_c2.sh}:

\begin{itemize}
    \item \textbf{Secuencial:} el código de la P1 compilado sin OpenMP. Es el que se usa como tiempo secuencial para calcular el speed-up.
    \item \textbf{P1 (OpenMP):} la mejor versión de la P1 con 2, 4, 6 y 12 hilos.
    \item \textbf{P2 (MPI):} la versión elegida en la cuestión 6 de la P2 (colectivas) con 2, 4, 6 y 12 procesos.
    \item \textbf{Híbrida:} todas las combinaciones de procesos e hilos con EP = 1, 2, 4, 6 y 12 elementos de proceso (EP = n\_procesos $\cdot$ n\_hilos).
\end{itemize}

En la copia del código de la P1 se han corregido dos errores que ya se habían corregido en la P2 para que la comparación sea justa: la función de ordenación \texttt{comp\_fitness} devolvía como \texttt{int} la resta de dos \texttt{double}, y en \texttt{distancia\_ij} estaba comentado el intercambio de \texttt{i} y \texttt{j}.

% >>> DATO: instancia, NGM, NEM y repeticiones
Se ha usado la instancia \texttt{n = 2000, m = 800, n\_gen = 400, tam\_pob = 200, m\_rate = 0.3}, con NGM = 10 y NEM = 3 (los valores elegidos en la P2), y cada configuración se ha ejecutado 3 veces tomando la media. Como en la P2, cada versión se ha ejecutado de dos formas:

\begin{itemize}
    \item \textbf{Trabajo fijo:} siempre 60 generaciones, sin criterio de convergencia (compilando con \texttt{-DSIN\_CONVERGENCIA}). Todas las versiones hacen el mismo trabajo, así que es la forma justa de comparar tiempos.
    \item \textbf{Con convergencia:} el algoritmo normal, que para cuando el fitness deja de mejorar. Con esta se compara tiempo y fitness con el indicador $I = \text{fitness}/\text{tiempo}$ de la observación 1.
\end{itemize}

% >>> DATO: tiempos de trabajo fijo del híbrido (resultados_c2/resumen.txt, filas "fijo hibrido")
\begin{table}[htbp]
\centering
\begin{tabular}{ccc|ccc}
\toprule
\textbf{n\_procesos} & \textbf{n\_hilos} & \textbf{EP} & \textbf{t\_ejec (s)} & \textbf{speed-up} & \textbf{eficiencia} \\
\midrule
1 & 1 & 1 & 14.803 & 1.01 & 1.01 \\
\midrule
1 & 2 & \multirow{2}{*}{2} & \textbf{8.333} & \textbf{1.79} & 0.89 \\
2 & 1 & & 11.239 & 1.33 & 0.66 \\
\midrule
1 & 4 & \multirow{3}{*}{4} & \textbf{6.342} & \textbf{2.35} & 0.59 \\
2 & 2 & & 8.298 & 1.80 & 0.45 \\
4 & 1 & & 10.304 & 1.45 & 0.36 \\
\midrule
1 & 6 & \multirow{4}{*}{6} & \textbf{6.054} & \textbf{2.46} & 0.41 \\
2 & 3 & & 8.053 & 1.85 & 0.31 \\
3 & 2 & & 8.408 & 1.77 & 0.30 \\
6 & 1 & & 9.752 & 1.53 & 0.25 \\
\midrule
1 & 12 & \multirow{6}{*}{12} & \textbf{4.563} & \textbf{3.26} & 0.27 \\
2 & 6 & & 5.978 & 2.49 & 0.21 \\
3 & 4 & & 6.959 & 2.14 & 0.18 \\
4 & 3 & & 7.610 & 1.96 & 0.16 \\
6 & 2 & & 8.240 & 1.81 & 0.15 \\
12 & 1 & & 9.980 & 1.49 & 0.12 \\
\bottomrule
\end{tabular}
\caption{Prestaciones de la versión híbrida con trabajo fijo (\texttt{n = 2000, m = 800, n\_gen = 60, tam\_pob = 200, m\_rate = 0.3}, NGM = 10, NEM = 3). El speed-up se calcula respecto al secuencial (14.896 s) y la eficiencia es speed-up/EP. En negrita, la mejor configuración para cada EP.}
\label{tab:c2_hibrido}
\end{table}

% >>> DATO: comparación de versiones (filas "fijo" y "convergencia" de resultados_c2/resumen.txt).
% Para el híbrido se pone la mejor configuración de cada EP.
\begin{table}[htbp]
\centering
\small
\begin{tabular}{lcc|cc|ccc}
\toprule
 & & & \multicolumn{2}{c|}{\textbf{Trabajo fijo}} & \multicolumn{3}{c}{\textbf{Con convergencia}} \\
\textbf{Versión} & \textbf{p $\times$ h} & \textbf{EP} & \textbf{t (s)} & \textbf{speed-up} & \textbf{t (s)} & \textbf{fitness} & \textbf{I ($\cdot 10^6$)} \\
\midrule
Secuencial & 1 $\times$ 1 & 1 & 14.896 & 1.00 & 3.975 & 16\,044\,825 & 4.04 \\
\midrule
P1 (OpenMP) & 1 $\times$ 2 & 2 & 7.882 & 1.89 & 2.212 & 16\,053\,970 & 7.26 \\
P2 (MPI) & 2 $\times$ 1 & 2 & 10.933 & 1.36 & 3.883 & 16\,048\,228 & 4.13 \\
Híbrida & 1 $\times$ 2 & 2 & 8.333 & 1.79 & 2.826 & 16\,048\,430 & 5.68 \\
\midrule
P1 (OpenMP) & 1 $\times$ 4 & 4 & 5.968 & 2.50 & 1.733 & 16\,045\,016 & 9.26 \\
P2 (MPI) & 4 $\times$ 1 & 4 & 10.980 & 1.36 & 3.680 & 16\,048\,737 & 4.36 \\
Híbrida & 1 $\times$ 4 & 4 & 6.342 & 2.35 & 2.172 & 16\,044\,370 & 7.39 \\
\midrule
P1 (OpenMP) & 1 $\times$ 6 & 6 & 5.322 & 2.80 & 1.585 & 16\,049\,307 & 10.13 \\
P2 (MPI) & 6 $\times$ 1 & 6 & 10.065 & 1.48 & 3.399 & 16\,048\,195 & 4.72 \\
Híbrida & 1 $\times$ 6 & 6 & 6.054 & 2.46 & 1.927 & 16\,045\,675 & 8.33 \\
\midrule
P1 (OpenMP) & 1 $\times$ 12 & 12 & 4.622 & 3.22 & \textbf{1.343} & 16\,044\,395 & \textbf{11.95} \\
P2 (MPI) & 12 $\times$ 1 & 12 & 9.724 & 1.53 & 3.395 & 16\,052\,520 & 4.73 \\
Híbrida & 1 $\times$ 12 & 12 & \textbf{4.563} & \textbf{3.26} & 1.614 & 16\,044\,175 & 9.94 \\
\bottomrule
\end{tabular}
\caption{Comparación de la versión secuencial, la mejor versión de la P1, la mejor de la P2 y la híbrida en el mismo sistema (\texttt{n = 2000, m = 800, tam\_pob = 200, m\_rate = 0.3}, NGM = 10, NEM = 3; 60 generaciones con trabajo fijo y hasta 400 con convergencia). $I = \text{fitness}/\text{tiempo}$.}
\label{tab:c2_comparacion}
\end{table}

% >>> DATO: los mismos tiempos de trabajo fijo de las tablas anteriores
\begin{figure}[htbp]
\centering
\begin{tikzpicture}
\begin{axis}[
    width=0.48\textwidth, height=6cm,
    xlabel={EP (procesos $\cdot$ hilos)}, ylabel={Tiempo (s)},
    xtick={1,2,4,6,12}, ymin=0,
    legend style={font=\footnotesize, at={(0.98,0.98)}, anchor=north east},
    grid=major, title={\small Tiempo con trabajo fijo}]
\addplot[mark=*, blue] coordinates {(1,14.896) (2,7.882) (4,5.968) (6,5.322) (12,4.622)};
\addlegendentry{P1 (OpenMP)}
\addplot[mark=square*, red] coordinates {(1,14.896) (2,10.933) (4,10.980) (6,10.065) (12,9.724)};
\addlegendentry{P2 (MPI)}
\addplot[mark=triangle*, teal] coordinates {(1,14.803) (2,8.333) (4,6.342) (6,6.054) (12,4.563)};
\addlegendentry{Híbrida 1 $\times$ h}
\end{axis}
\end{tikzpicture}
\hfill
\begin{tikzpicture}
\begin{axis}[
    width=0.48\textwidth, height=6cm,
    xlabel={n\_procesos (n\_hilos = 12/n\_procesos)}, ylabel={Tiempo (s)},
    xtick={1,2,3,4,6,12}, ymin=0,
    grid=major, title={\small Híbrida con EP = 12}]
\addplot[mark=triangle*, teal] coordinates {(1,4.563) (2,5.978) (3,6.959) (4,7.610) (6,8.240) (12,9.980)};
\end{axis}
\end{tikzpicture}
\caption{Izquierda: tiempo con trabajo fijo de cada versión al aumentar los elementos de proceso. Derecha: tiempo de la versión híbrida con 12 elementos de proceso repartidos de distintas formas entre procesos e hilos.}
\label{fig:c2}
\end{figure}

\medskip
\noindent\textbf{Análisis}\par
\smallskip

% >>> DATO: revisar los números citados si cambian las tablas
\begin{itemize}
    \item \textbf{Con el mismo número de elementos de proceso, es mejor usar más hilos y menos procesos.} En la Tabla~\ref{tab:c2_hibrido}, para cada EP la mejor configuración es siempre la de 1 proceso, y el tiempo empeora a medida que se pasan hilos a procesos. Con EP = 12, por ejemplo, se pasa de 4.56 s con 1 $\times$ 12 a 9.98 s con 12 $\times$ 1 (Figura~\ref{fig:c2}, derecha). Hay dos motivos. El primero es la memoria: cada proceso MPI tiene su propia copia de la matriz de distancias (unos 16 MB con $n = 2000$), mientras que los hilos de un proceso comparten una sola. Con 12 procesos hay 12 copias que no caben en la caché y compiten por el acceso a memoria, y el fitness, que es la parte más costosa, se pasa el tiempo leyendo esa matriz. El segundo son las migraciones: cada NGM generaciones todos los procesos tienen que esperar al más lento, y con más procesos esa espera es mayor.
    \item \textbf{Respecto al secuencial.} La mejor configuración híbrida es 1 proceso con 12 hilos, con un speed-up de 3.26 y una eficiencia de 0.27. La eficiencia baja tanto por las características del procesador: solo 2 de los 10 núcleos son de rendimiento, el resto son más lentos, y con 12 hilos dos pares de hilos comparten núcleo. Además, en un portátil la frecuencia baja cuando se usan todos los núcleos a la vez. Con 1 proceso y 1 hilo el tiempo es igual que el secuencial, así que la versión híbrida casi no añade coste cuando no se usa el paralelismo.
    \item \textbf{Respecto a la mejor versión de la P1.} Con 1 proceso, la versión híbrida es prácticamente la de la P1 más las migraciones, que con una sola isla no aportan nada. Por eso es entre un 6 y un 14\,\% más lenta con 2 a 6 hilos, y con 12 hilos tardan lo mismo (4.56 s frente a 4.62 s, una diferencia que entra en la variación entre ejecuciones). En un único ordenador la memoria compartida es lo más eficiente, y el nivel MPI solo compensaría con varios nodos, donde OpenMP no llega por sí solo.
    \item \textbf{Respecto a la mejor versión de la P2.} La versión híbrida es mucho mejor: con EP = 12 tarda 4.56 s frente a 9.72 s, es decir, 2.1 veces más rápida. La versión de la P2 no pasa de un speed-up de 1.53, por los mismos motivos de memoria explicados antes.
    \item \textbf{Tiempo y fitness.} El fitness es prácticamente igual en todas las versiones y configuraciones (entre 16.04 y 16.07 millones, menos de un 0.2\,\% de diferencia). Por eso el indicador $I$ ordena las versiones casi igual que el tiempo. Con convergencia, la mejor versión híbrida (1 $\times$ 12) consigue $I = 9.94 \cdot 10^6$, 2.5 veces el del secuencial y 2.1 veces el de la mejor de la P2, aunque algo por debajo de la P1 con 12 hilos ($11.95 \cdot 10^6$).
    \item \textbf{Diferencia con convergencia.} Con convergencia la versión híbrida sale peor que la P1, y con 1 proceso y 1 hilo incluso más lenta que el secuencial (5.14 s frente a 3.98 s). No es porque sea menos eficiente: la P1 comprueba la convergencia en cada generación, mientras que la versión con islas solo puede hacerlo en cada migración, cada NGM = 10 generaciones, y por tanto hace más generaciones antes de parar. Por eso, para comparar el rendimiento de las implementaciones, la medida fiable es la de trabajo fijo, y la de convergencia sirve para ver el compromiso entre tiempo y fitness.
\end{itemize}

En resumen, en este sistema la mejor configuración de la versión híbrida es 1 proceso con 12 hilos, que es más de 3 veces más rápida que el secuencial, el doble de rápida que la mejor versión de la P2 y tan rápida como la mejor de la P1.


% =====================================================================
%  PARTE 2
% =====================================================================
\newpage
\item[] \textbf{\large Parte 2. Modelado y auto-optimización del algoritmo paralelo}

En toda la Parte 2 se ha sustituido el fichero \texttt{io.c} por el proporcionado en \texttt{datos/ejecucion/}, que genera la matriz de distancias directamente en memoria, sin ficheros de entrada.

\item \textbf{Cuestión 3.} Modelo de tiempo en memoria compartida con 1 nivel de paralelismo.

% --- Respuesta de la cuestión 3 ---

\medskip
\noindent\textbf{3.1 Obtención de datos experimentales}\par
\smallskip

% a) paralelizar FOR_INI (con el fitness dentro)
% b) nuevos parámetros N_HILOS_INI y N_HILOS_FIT
% c) medir el tiempo de FOR_INI con h = 1, 2, 4, ..., 20
% d) n = 30000, m = 28000, NE = 6, n_gen = 0, 5 series

\medskip
\noindent\textbf{3.2 Instalación de la rutina con 1 nivel de paralelismo}\par
\smallskip

% ajuste de ks1 y kh en la hoja "Modelo 1-Nivel Paralelismo" de solver.ods

\medskip
\noindent\textbf{3.3 Ejecución}\par
\smallskip

% a) hopt = sqrt(ks1 * NE / kh)
% b) y c) predicción y comprobación con h = 3, 5, 7, ...
% d) y e) hopt con NE = 3 y NE = 10, teórico frente a experimental


\newpage
\item \textbf{Cuestión 4.} Modelo de tiempo en memoria compartida con 2 niveles de paralelismo.

% --- Respuesta de la cuestión 4 ---

\medskip
\noindent\textbf{4.1 Instalación de la rutina con 2 niveles de paralelismo}\par
\smallskip

% paralelizar FOR_INI y FOR_FIT con omp_set_nested(1); ajuste de ks2, kh1 y kh2

\medskip
\noindent\textbf{4.2 Número óptimo de hilos}\par
\smallskip

% h1opt y h2opt a partir de las ecuaciones 4 y 5

\medskip
\noindent\textbf{4.3 Ejecución}\par
\smallskip

% a) tiempo esperado con otros valores de h1 y h2
% b) h1opt y h2opt con otros valores de NE y |B|


\newpage
\item \textbf{Cuestión 5.} Modelo de tiempo en memoria distribuida.

% --- Respuesta de la cuestión 5 ---

\medskip
\noindent\textbf{5.1 Tiempo de ordenación}\par
\smallskip

% nueva hoja en solver.ods; comprobar que t_ord es despreciable

\medskip
\noindent\textbf{5.2 Instalación: constantes kcmp, A, B, C y D}\par
\smallskip

% NGM = 1, variando el número de procesos (Makefile de datos/ejecucion/pm)

\medskip
\noindent\textbf{5.3 Número óptimo de procesos}\par
\smallskip

% tiempo teórico para varios p y popt resolviendo la ecuación 10


\newpage
\item \textbf{Cuestión 6.} Realizar el informe final incluyendo las conclusiones generales y valoración personal sobre el trabajo realizado.

% --- Respuesta de la cuestión 6 ---

\end{enumerate}

% =========================
% COORDINACIÓN
% =========================
\newpage
\section{Coordinación, reparto del trabajo y tiempo dedicado}

% --- Rellenar al final ---

\end{document}